\scanpage{76}
\mspage{K}
% archive foliation: 75
\uline{Théorème} Soit $K$ une extension de type fini du
corps $k$, \deleted{[ayant au moins un modèle
propre \inserted{sur $k$} et universellement normal sur $k$]}
\deleted{Soit} Il
\deleted{$(X_i)_{i\in I}$ le système \unsure{ind}}
existe un entier $n(K/k)$
tel que pour tout modèle $X$ de $K$ propre sur $k$, et
normal sur $k$, \deleted{\unsure{pour}} on ait \deleted{$\dim\uline{\mathrm{Pic}}_X$}
\[
\dim\uline{\mathrm{Pic}}_{X/k}\leq n(K/k)
\]

\uline{Démonstration}. \deleted{\unsure{On est ramené au cas}}
\inserted{S'il n'existe pas de}
modèle $X$ du type indiqué, c'est trivial.
Sinon $K$ est séparable sur $k$, et
$\overline K=K\otimes_k\overline k$ est une extension de type
fini du corps alg. clos $\overline k$, et on
est ramené au cas où $k$ est alg. clos.
Alors
\[
\dim\uline{\mathrm{Pic}}_{X/k}
=\tfrac12\operatorname{rang}_{\mathbf Z_\ell}
(\text{partie $\ell$-primaire de }\pi_1(X/k)\text{ rendu abélien})
\]
et à ce titre, \unsure{on renvoie aux} théorèmes
de majoration uniforme des $\pi_1(X/k)$
